% Generated draft: requires mathematical and source review.
\documentclass[11pt]{article}
\usepackage[margin=25mm]{geometry}
\usepackage{fontspec}
\setmainfont{Latin Modern Roman}
\usepackage{amsmath,amssymb,mathtools,unicode-math}
\setmathfont{Latin Modern Math}
\usepackage{enumitem,needspace,xurl,hyperref}
\setlength{\parindent}{0pt}
\setlength{\parskip}{5pt}
\setlength{\emergencystretch}{3em}
\newcommand{\N}{\mathbb{N}}
\newcommand{\Z}{\mathbb{Z}}
\newcommand{\Q}{\mathbb{Q}}
\newcommand{\R}{\mathbb{R}}
\newcommand{\C}{\mathbb{C}}
\newcommand{\F}{\mathbb{F}}
\newcommand{\A}{\mathbb{A}}
\newcommand{\PP}{\mathbb{P}}
\newcommand{\E}{\mathbb{E}}
\newcommand{\eps}{\varepsilon}
\newcommand{\dd}{\,\mathrm{d}}
\begin{document}
\section*{UnsolvedMath ID: 20001864}
\textbf{Source problem code:} catalogue entry
\subsection{Definitions and mathematical statement}
Schulze's source asks when a linkage preserves its original symmetry throughout configuration space. The meaning of symmetry can differ: an automorphism of the length-labeled graph; invariance of one configuration modulo Euclidean motions; setwise preservation of a moduli component; pointwise fixing of every class in it; or realization by one fixed spatial representation in a chosen gauge. The aid's theorem addresses pointwise moduli symmetry on a connected smooth component, not all these meanings at once.

Let a finite length-labeled graph define its realization space, modulo a specified Euclidean-motion group. Let \(C\) be a connected smooth Riemannian moduli component, and let a length-preserving graph automorphism \(h\) act as a Riemannian isometry of \(C\). Suppose \(h(x)=x\) for one class \(x\in C\). The reported criterion is
\[
 h|_C=\operatorname{id}\quad\Longleftrightarrow\quad
 dh_x=\operatorname{id}_{T_xC}.
\]
The isometry hypothesis is essential; a general smooth diffeomorphism is not determined by its first derivative at a fixed point. The component and quotient regularity assumptions are also essential to this stated theorem.

At a regular full-dimensional framework \(p=(p_i)\), let \(\mathcal F_p\) be its infinitesimal-flex space, defined by \((p_i-p_j)\cdot(u_i-u_j)=0\) for every bar. Let
\[
 \mathcal T_p=\{u_i=a+A p_i:\ A^{\mathsf T}=-A\}
\]
be the trivial-motion space, so the intended smooth tangent quotient is \(\mathcal F_p/\mathcal T_p\). If \(\rho_h(p_i)=p_{h i}\) realizes the initial symmetry, with orthogonal linear part \(R_h\), the compensated action can be written \((L_hu)_i=R_h^{-1}u_{h i}\). The aid reports the equivalent check
\[
 L_hu-u\in\mathcal T_p\qquad\text{for every }u\in\mathcal F_p.
\]
For a symmetry group it suffices to check generators under the same fixed-class/isometric-action hypotheses. No such smooth one-point criterion is asserted through singular strata or on other disconnected components.

The aid's exact example is a unit four-bar rhombus family, in a pinned gauge representable by
\[
 p_0=0,\quad p_1=1,\quad p_2=1+e^{i\theta},\quad p_3=e^{i\theta},
 \qquad0<\theta<\pi.
\]
The quarter-cycle relabeling sends \(\theta\) to \(\pi-\theta\), preserving the component setwise and fixing the square at \(\theta=\pi/2\), but has derivative \(-1\) there and breaks pointwise symmetry immediately. Its square fixes every moduli class, realized by the half-turn \(z\mapsto1+e^{i\theta}-z\). That realizing Euclidean motion varies with \(\theta\) in this pinned gauge, despite pointwise moduli persistence; another gauge may change its spatial formula. This distinguishes a fixed moduli action from an unchanged spatial representation in prescribed coordinates. The example and criterion do not classify every linkage or guarantee symmetry across the full possibly singular configuration space.
\subsection{Short English statement}
On a connected smooth moduli component with an isometric symmetry action, the aid reduces pointwise persistence to the action on one infinitesimal-flex quotient. A unit four-bar distinguishes setwise preservation, a symmetry at one point, pointwise moduli persistence, and gauge-dependent spatial realizations.
\subsection{Sources}
\begin{itemize}
\item Local UnsolvedMath registry, permanent ID \(20001864\), source code \texttt{catalogue entry}; \texttt{aim-geometry-notes.json}, \texttt{notes[201]}, tag \texttt{problem}.
\item AIM workshop ``Configuration spaces of linkages,'' recorded item 5, attributed to Schulze: \url{https://aimath.org/pastworkshops/linkagesproblems.pdf}. Source type PDF; statement status \texttt{exact}. The supplied verification says the item number lost its first digit, but truncates before giving a corrected number; none is guessed. The trailing ``1'' is an extraction artifact. Only the supplied text was used; no PDF was accessed or compiled.
\item No authoritative bibliography, full quotient/metric conventions, or complete proof artifacts accompany the differential criterion or four-bar example. Original remarks and literature fields are empty.
\item Research attempt 1: \texttt{partial\_result}, status at run \texttt{open}, rigor label \texttt{proved}, novelty confidence low; registry status \texttt{partially\_solved}, classification \texttt{PARTIAL-PROGRESS}. Created and updated 14 August 2026. Machine-generated aid, not peer reviewed; independently verify before citation.
\item Ulam AI / UnsolvedMath Contributors; CC BY 4.0.
\end{itemize}
\end{document}
